\lecture{3}
\title{Kernel and Image, Cosets, Normal Subgroups, Conjugation, The Correspondence Theorem}

\begin{document}

\subsection{Kernel and Image}

\textbf{Definition (Kernel and Image).} Consider a homomorphism $f: G \to H$ between groups $G$ and $H$.  Then the \textbf{kernel} of $f$ is $\ker(f) := \{g \in G : f(g) = 1_H\}$, and the \textbf{image} of $f$ is $\im (f) := \{h \in H: h = f(g) \text{ for some } g \in G\}$.

Note that $\ker(f)$ is a subgroup of $G$, and $\im(f)$ is a subgroup of $H$.  (Check this!)

\begin{problem}{Application of the Kernel}
    Take a homomorphism $f: G \to H$, and suppose that $f(g_1) = f(g_2)$ for $g_1, g_2 \in G$.  What can we say about $g_1$ and $g_2$?
\end{problem}

\begin{solution*}
    Observe the following:
    \[ f(g_1) = f(g_2) ~ \iff ~ f(g_1g_2^{-1}) = 1 ~ \iff ~ g_1g_2^{-1} \in \ker(f) ~ \iff ~ g_1 = kg_2 ~ \text{ for some } k \in \ker(f). \]
    If we multiply by the left inverse instead of the right inverse in the first step, we also find:
    \[ f(g_1) = f(g_2) ~ \iff ~ f(g_2^{-1}g_1) = 1 ~ \iff ~ g_2^{-1}g_1 \in \ker(f) ~ \iff ~ g_1 = g_2k \text{ for some } k \in \ker(f). \]
    Consolidating these results together, we have:
    \[ \boxed{\begin{cases}f(g_1) = f(g_2) ~ \iff ~ g_1 = kg_2 \text{ for some } k \in \ker(f) \\
    f(g_1) = f(g_2) ~ \iff ~ g_1 = g_2k \text{ for some } k \in \ker(f)\end{cases}} \]
\end{solution*}

\textbf{Example.} Take the homomorphism $f: (\mathbb{C}, +) \to (\mathbb{C}^{\times}, \cdot)$ defined by $f(z) := e^z$.  Then the image of $f$ is $\im(f) = \mathbb{C}^{\times}$ (i.e. $f$ is surjective), and the kernel of $f$ is $\ker(f) = 2\pi i \cdot \mathbb{Z}$.

\textbf{Example.} Take the homomorphism $f: GL_n(\mathbb{R}) \to \mathbb{R}^{\times}$ defined by $f(M) = \det M$.  Then the image of $f$ is $\im(f) = \mathbb{R}^{\times}$ (i.e. $f$ is surjective), and the kernel of $f$ is $\ker(f) = SL_n(\mathbb{R})$.

\textbf{Example.} Take the homomorphism $f: S_n \to \{-1, 1\}$ send permutations $\pi \in S_n$ to the sign of $\pi$.  The kernel of $f$ must be a subgroup of $S_n$; this subgroup is the \textbf{alternating group} $A_n$, consisting of all even permutations in $S_n$.

\textbf{Example.} Let $G$ be a group, and fix an element $g \in G$.  Then we may take the homomorphism $f: (\mathbb{Z}, +) \to G$ defined by $f(n) := g^n$.  In this case,
\begin{itemize}
    \item The image of $f$ is the subgroup generated by $\{g\}$.  We denote this by $\im(f) = \langle g \rangle$.
    \item The kernel of $f$ is $\ker(f) = k\mathbb{Z}$, where $g$ has order $k$.  (If $g$ has infinite order, then $\ker(f) = \{0\}$.)
\end{itemize}

\subsection{Cosets}

\textbf{Example.} Recall that $3\mathbb{Z}$ is a subgroup of $\mathbb{Z}$.  This subgroup can be ``copied-and-translated'' twice to form the subsets $3\mathbb{Z} + 1$ and $3\mathbb{Z} + 2$ in $G$.  (Note that these subsets are not subgroups!)

Let's give a name to these ``copied-and-translated'' subsets associated with a subgroup: \textbf{cosets.}

\textbf{Definition (Cosets).} Let $H$ be a subgroup of $G$.  Then a \textbf{left coset} of $H$ is a subset $gH := \{gh : h \in H\} \subseteq G$ for some $g \in G$.  Similarly, a \textbf{right coset} of $H$ is a subset $Hg := \{hg : h \in H\} \subseteq G$ for some $g \in G$.

The phrase ``coset'' (without further specification) should be assumed to mean ``left coset''.

\begin{theorem}{Coset Properties}
    Let $H$ be a subgroup of $G$.  The cosets of $H$ in $G$ satisfy the following properties:
    \begin{enumerate}
        \item Any two cosets of $H$ have the same size.
        \item Any two distinct cosets are disjoint.
    \end{enumerate}
\end{theorem}

\begin{proof}
    To prove the first fact, simply note that the size of any coset $gH$ is equal to $|H|$.  This is because the natural map $T: H \to gH$ defined by $T(h) = gh$ is a bijection. ~ $\blacksquare$

    To prove the second fact, consider two distinct cosets $C_1 = g_1H$ and $C_2 = g_2H$.  Suppose, for the sake of contradiction, that these two cosets are not disjoint.  Then there exists $h_1, h_2 \in H$ such that $g_1h_1 \in C_1$ and $g_2h_2 \in C_2$ are the same; that is, $g_1h_1 = g_2h_2$.  Equivalently, $g_1 = g_2h_2h_1^{-1}$.  However, this implies
    \[ C_ 1 = g_1H = (g_2h_2h_1^{-1})H = g_2(h_2h_1^{-1}H) = g_2H = C_2. \]
    So in fact, we have shown that if there is just a \emph{single} pair of elements $g_1h_1$ and $g_2h_2$ in $C_1$ and $C_2$ that match, then we can show the \emph{entirety} of $C_1$ and $C_2$ must match completely!  ~ $\blacksquare$
\end{proof}

Here's an equivalent, more intuitive rephrasing of the above theorem:

\begin{theorem}{Cosets form Partitions}
    Let $H$ be a subgroup of $G$.  Then the set of all left cosets of $H$ partition the elements of $G$ into disjoint subsets of size $|H|$.
\end{theorem}

Note that all of the above theorems hold equally true if every instance of ``left'' is replaced with ``right''.  However, the partitioning of $G$ formed by left cosets usually does not equal the partitioning formed by right cosets!

\textbf{Example.} Consider the group $G = (\mathbb{R}^2, +)$ and a subgroup $H := \{(a, 2a): a \in \mathbb{R}\}$.  Then cosets of $H$ look like translated copies of lines parallel to $y = 2x$.  These lines tile the plane $\mathbb{R}^2$.

\textbf{Definition (Index).} The number of cosets of $H$ in $G$ is the \textbf{index} of $H$ in $G$, denoted $[G:H]$.  Note that the theorem above implies $|G| = [G:H] \cdot |H|$.

\begin{theorem}{Lagrange's Theorem}
    Let $H$ be a (finite) subgroup of a (finite) group $G$.  Then $|H|$ divides $|G|$.
\end{theorem}

\begin{proof}
    Follows from the equation $|G| = [G:H] \cdot |H|$, along with the fact that $[G:H]$ is an integer. ~ $\blacksquare$
\end{proof}

Lagrange's Theorem is very powerful!  Here's an interesting corollary of the above theorem.

\begin{theorem}{Groups of Prime Order}
    Any group of prime order must be cyclic.
\end{theorem}

\begin{proof}
    Suppose $G$ is a group of prime order $p$.  Pick some $g \in G$ that is not the identity.  Then $H = \langle g \rangle$ is a subgroup of $G$.  Importantly, $|H| > 1$ (since $g$ is not the identity), and $|H|$ divides $p$.  So it must be that $|H| = p$, meaning that $G$ and $H$ are the same group.  And since $H$ was a cyclic group, $G$ must be cyclic, too.
\end{proof}

\subsection{Normal Subgroups}

How do cosets and homomorphisms relate?

\begin{problem}{Application of the Coset}
    Take a homomorphism $f: G \to H$, and suppose that $f(g_1) = f(g_2)$ for $g_1, g_2 \in G$.  What can we say about $g_1$ and $g_2$?
\end{problem}

\begin{solution*}
    Recall the following result we derived earlier:
    \[ \boxed{\begin{cases}f(g_1) = f(g_2) ~ \iff ~ g_1 = kg_2 \text{ for some } k \in \ker(f) \\
    f(g_1) = f(g_2) ~ \iff ~ g_1 = g_2k \text{ for some } k \in \ker(f)\end{cases}} \]
    We can now use cosets to rephrase the above statement as follows:
    \[ \boxed{\begin{cases}f(g_1) = f(g_2) ~ \iff ~ g_1 \text{ and } g_2 \text{ are in the same left coset of } \ker(f). \\
    f(g_1) = f(g_2) ~ \iff ~ g_1 \text{ and } g_2 \text{ are in the same right coset of } \ker(f). \end{cases}} ~~~~~ (\bigstar) \]
\end{solution*}

The above motivates the following definition:

\textbf{Definition (Normal).} A subgroup $N \subseteq G$ of $G$ is \textbf{normal} if $gN = Ng$ for all $g \in G$.

This is because the above result $(\bigstar)$ immediately implies the following theorem:

\begin{theorem}{Normal Kernel}
    For a homomorphism $f: G \to H$, its kernel $K = \ker(f)$ is a normal subgroup of $G$.
\end{theorem}

\textbf{Example.} Consider $G = S_3$ and a subgroup $N = \{e, (1 \ 2)\}$.  Then:
\begin{itemize}
    \item The left cosets of $N$ are $N = \{e, (1 \ 2)\}$, $N(1 \ 3) = \{(1 \ 3), (1 \ 3 \ 2)\}$, and $N(2 \ 3) = \{(2 \ 3), (1 \ 2 \ 3)\}$.
    \item The right cosets of $N$ are $N = \{e, (1 \ 2)\}$, $(1 \ 3)N = \{(1 \ 3), (1 \ 2 \ 3)\}$, and $(2, 3)N = \{ (2 \ 3), (1 \ 3 \ 2)\}$.
\end{itemize}
Thus, $N$ is not normal, since (e.g.) $N(1 \ 3) \neq (1 \ 3)N$.  In other words, the partitioning of $G$ formed by the left cosets of $N$ is \emph{different} from the partitioning of $G$ formed by the right cosets of $N$.  

\begin{remark}
    By corollary, there is no homomorphism $f: G \to H$ to \emph{any} other group $H$ such that $\ker(f) = N$.
\end{remark}

\textbf{Example.} If $G$ is abelian, then every subgroup $N \subseteq G$ is normal.  This is because:
\[ gN = \{gn : n \in N\} = \{ng : n \in N\} = Ng. \]

\textbf{Example.} If $N$ is a subgroup of $G$ such that $[G:N] = 2$, then $N$ must be normal.  This is because any partitioning of $G$ into cosets---whether they be left cosets or right cosets---must look like $G = N \  \sqcup \ (G \setminus N)$.  So the (two) left cosets of $N$ must match the (two) right cosets of $N$.

\subsection{Conjugation}

\textbf{Definition (Conjugation).} Given elements $a, b \in G$, the \textbf{conjugation} of $a$ by $b$ is defined to be $bab^{-1}$.  More generally, if $H$ is a subgroup of $G$, then the \textbf{conjugation} of $a$ by $H$ is $aHa^{-1}$.

It turns out (easy to check) that if $H$ is a subgroup, then $aHa^{-1}$ is also always a subgroup.  Thus, we can think of conjugation by $g \in G$ as an action that sends subgroups of $G$ to subgroups of $G$.

\begin{theorem}{Normality by Conjugation Condition}
    A subgroup $N \subseteq G$ is normal if and only if $N$ does not change under conjugation by any element $g \in G$.
\end{theorem}

\begin{proof}
    Recall that $N$ is normal if and only if $gN = Ng$ for all $g \in G$.  Equivalently, $gNg^{-1} = N$ for all $g \in G$, which may directly be read as: the conjugation of $N$ by $g$ equals $N$ for any $g \in G$.  ~ $\blacksquare$
\end{proof}

\begin{theorem}{Intersecting Normal Subgroups}
    If $N_1$ and $N_2$ are normal subgroups of $G$, then their intersection $N_1 \cap N_2$ is also normal.
\end{theorem}

\begin{proof}
    We argue using the formulation of normality in terms of conjugation described above.  For any $g \in G$,
    \begin{align*} g(N_1 \cap N_2)g^{-1} = \ & \{gng^{-1} : n \in N_1 \cap N_2\} \\ = \ & \{gng^{-1} : n \in N_1\} \cap \{gng^{-1}: n \in N_2\} \\
    = \ & (gN_1g^{-1}) \cap (gN_2g^{-1}) \\
    = \ & N_1 \cap N_2 \text{ since } N_1 \text{ and } N_2 \text{ are normal.} \end{align*}
    Thus, $N_1 \cap N_2$ remains unchanged by conjugation by $g$, so $N_1 \cap N_2$ is normal. ~ $\blacksquare$
\end{proof}

\subsection{Correspondence Theorem}

Here's a useful theorem that we won't prove today.

\begin{theorem}{Correspondence Theorem}
    Let $f: G \to G'$ be a \emph{surjective} homomorphism, and let $K = \ker(f)$ be its kernel.  Consider the following collections of subgroups:
    \[ A := \{\text{subgroups of } G \text{ containing } K\} ~ \text{ and } ~ B := \{\text{subgroups of } G'\}. \]
    Then there is a bijection from $A$ to $B$ that sends subgroups $H \in A$ to subgroups $f(H) := \{f(h): h \in H\} \in B$.
\end{theorem}

\begin{problem}{Correspondence Theorem Application \#1}
    Consider the determinant homomorphism $\det: G \to G'$, with $G = GL_n(\mathbb{R})$ and $G' = \mathbb{R}^{\times}$.  Then $K = \ker(\det) = SL_2(\mathbb{R})$.  What does the Correspondence Theorem say?
\end{problem}

\begin{solution}
    The set $B := \{\text{subgroups of } \mathbb{R}^{\times}\}$ is easier to analyze; for example, $\mathbb{Q}^{\times}$ is a subgroup in $B$.
    
    By the Correspondence Theorem, the subgroup $\mathbb{Q}^{\times}$ must correspond to a new, quirky subgroup of $GL_n(\mathbb{R})$ containing $K = SL_2(\mathbb{R})$ that we didn't know of before.  This subgroup is:
    \[ f^{-1}(\mathbb{Q}^{\times}) = \{M \in GL_n(\mathbb{R}): \det M \in \mathbb{Q}^+\} \in A. \]
    There are many, many other subgroups of $GL_n(\mathbb{R})$ that live between $K$ and $G$.
\end{solution}

\begin{problem}{Correspondence Theorem Application \#2}
    Take $G = \mathbb{Z}$ and $G' = C_n = \langle g \rangle$, and define the homomorphism $f: G \to G'$ by $f(k) = g^k$.  What does the Correspondence Theorem say?
\end{problem}

\begin{solution}
    Observe that $\ker(f) = n \mathbb{Z}$.  Then:
    \[ A := \{\text{subgroups of } G \text{ containing } K\} = \{d\mathbb{Z} : d \mid n\}. \]
    It then follows that we can classify the set of all subgroups of $G' = C_n$:
    \[ B := \{ \text{subgroups of } G'\} = f(A) = \{f(d\mathbb{Z}) : d \mid n\} = \{ \langle g^d \rangle : d \mid n \}. \]
    So in fact, this tells us every subgroup of $C_n$ is of the form $\langle g^d \rangle$ for some $d \mid n$.
\end{solution}

% test2
% test3
